\documentclass[11pt,a4paper]{article}

\usepackage[utf8]{inputenc}
\usepackage[T1]{fontenc}
\usepackage{lmodern}
\usepackage[spanish]{babel}
\usepackage[margin=2cm]{geometry}
\usepackage{microtype}
\usepackage{amsmath}
\usepackage{amssymb}
\usepackage{xcolor}
\usepackage{hyperref}
\usepackage{booktabs}
\usepackage{longtable}
\usepackage{array}
\usepackage{tabularx}
\usepackage{multirow}
\usepackage{enumitem}
\usepackage{listings}
\usepackage{framed}
\usepackage{titlesec}
\usepackage{parskip}
\usepackage{tikz}
\usetikzlibrary{shapes.geometric, arrows.meta, positioning, fit, backgrounds, calc}
\usepackage{rotating}
\usepackage{pdflscape}

\hypersetup{
  colorlinks=true,
  linkcolor=blue!50!black,
  urlcolor=blue!60!black,
  pdftitle={Pipeline de datos a modelo --- POC ValueData},
  pdfauthor={ValueData}
}

% --- Colores ---
\definecolor{stage1}{RGB}{220,235,250}    % Excel
\definecolor{stage2}{RGB}{225,245,225}    % SQLite
\definecolor{stage3}{RGB}{255,240,210}    % Generador sintetico
\definecolor{stage4}{RGB}{245,225,250}    % Feature engineering
\definecolor{stage5}{RGB}{250,225,225}    % Train
\definecolor{stage6}{RGB}{225,235,235}    % Inferencia
\definecolor{codebg}{RGB}{248,248,248}
\definecolor{codekw}{RGB}{0,80,160}
\definecolor{codecmt}{RGB}{120,130,120}
\definecolor{codestr}{RGB}{160,80,40}
\definecolor{accent}{RGB}{0,100,180}

% --- Listings ---
\lstdefinestyle{mono}{
  basicstyle=\ttfamily\small,
  backgroundcolor=\color{codebg},
  keywordstyle=\color{codekw}\bfseries,
  commentstyle=\color{codecmt}\itshape,
  stringstyle=\color{codestr},
  breaklines=true, breakatwhitespace=true, showstringspaces=false,
  frame=single, framesep=4pt, framerule=0pt, rulecolor=\color{codebg},
  xleftmargin=8pt, xrightmargin=8pt, aboveskip=8pt, belowskip=8pt,
  literate={á}{{\'a}}1 {é}{{\'e}}1 {í}{{\'i}}1 {ó}{{\'o}}1 {ú}{{\'u}}1
           {Á}{{\'A}}1 {É}{{\'E}}1 {Í}{{\'I}}1 {Ó}{{\'O}}1 {Ú}{{\'U}}1
           {ñ}{{\~n}}1 {Ñ}{{\~N}}1 {¿}{{?`}}1 {¡}{{!`}}1
           {≈}{{$\approx$}}1 {×}{{$\times$}}1 {≥}{{$\geq$}}1 {≤}{{$\leq$}}1
           {→}{{$\rightarrow$}}1 {∧}{{$\wedge$}}1 {…}{{\ldots}}1
           {°}{{$^{\circ}$}}1 {²}{{$^{2}$}}1 {³}{{$^{3}$}}1
}
\lstset{style=mono}
\lstdefinelanguage{pyplus}{
  language=Python,
  morekeywords={async,await,with,as,from,import,return,for,if,else,elif,def,class,not,and,or,True,False,None,in,is,raise,try,except,finally,lambda},
}

% --- Bloque "stage" estilo card ---
\newenvironment{stagebox}[2]{%
  \par\smallskip
  \noindent\begin{tikzpicture}
    \node[draw=accent, thick, rounded corners=2pt, fill=#1, inner sep=8pt,
          text width=\dimexpr\linewidth-22pt\relax, align=left] (b) {%
      \textcolor{accent}{\textbf{#2}}\par\smallskip\ignorespaces
}{%
    };
  \end{tikzpicture}\par\smallskip
}

\definecolor{shadecolor}{RGB}{240,247,255}
\newenvironment{quoteblock}{\begin{shaded}\noindent\ignorespaces}{\end{shaded}}

\titleformat{\section}[block]{\Large\bfseries\color{accent}}{\thesection.}{0.6em}{}
\titleformat{\subsection}{\large\bfseries}{\thesubsection}{0.5em}{}
\titlespacing*{\section}{0pt}{1.4em}{0.7em}
\titlespacing*{\subsection}{0pt}{0.9em}{0.4em}

\setlist[itemize]{leftmargin=1.4em,itemsep=2pt,topsep=4pt}
\setlist[enumerate]{leftmargin=1.6em,itemsep=2pt,topsep=4pt}

\title{\textbf{Pipeline de datos a modelo}\\[0.3em]
       \large POC ValueData --- de las columnas del Excel al clasificador en producción}
\author{}
\date{}

\begin{document}
\maketitle
\thispagestyle{empty}
\vspace{-2em}

\begin{quoteblock}
\textbf{Objetivo de este documento}: trazar paso a paso el viaje de los datos desde el Excel que entregó Falabella (\texttt{datos\_eta\_*.xlsx}) hasta las predicciones que el frontend muestra en la torre. Cada \textbf{stage} es una unidad reproducible con \emph{input}, \emph{transformación}, \emph{output} y \emph{código de referencia}. Al final hay un linaje columna por columna (Apéndice A) y un checklist de validación (Apéndice B).
\end{quoteblock}

% =====================================================================
\section{Vista general del pipeline}
% =====================================================================

\begin{center}
\begin{tikzpicture}[
  node distance=4mm and 8mm,
  every node/.style={font=\small},
  stage/.style={rectangle, rounded corners=2pt, draw=accent, thick,
                minimum width=3.4cm, minimum height=1.05cm,
                align=center, inner sep=4pt},
  arrow/.style={-Stealth, thick, draw=accent!70!black},
  branch/.style={-Stealth, thick, draw=gray, dashed}
]
% Track A: analytics
\node[stage, fill=stage1] (s1) {\textbf{1.}~Excel\\\footnotesize Simpli + Geo};
\node[stage, fill=stage2, right=of s1] (s2) {\textbf{2.}~SQLite\\\footnotesize \texttt{fpoc\_*}};
\node[stage, fill=stage6, right=of s2] (s2b) {\textbf{2b.}~Seguimiento\\\footnotesize KPIs, motivos};

% Track B: predictivo
\node[stage, fill=stage3, below=12mm of s1] (s3) {\textbf{3.}~Generador\\\footnotesize sintético 60d};
\node[stage, fill=stage4, right=of s3] (s4) {\textbf{4.}~Feature\\engineering};
\node[stage, fill=stage4, right=of s4] (s5) {\textbf{5.}~Split\\\footnotesize 50/10 días};
\node[stage, fill=stage5, below=12mm of s5] (s6) {\textbf{6.}~Train\\\footnotesize XGB + isotonic};
\node[stage, fill=stage5, left=of s6] (s7) {\textbf{7.}~Pruebas\\\footnotesize AUC, Brier, calib};
\node[stage, fill=stage6, left=of s7] (s8) {\textbf{8.}~Inferencia\\\footnotesize cada 3s};
\node[stage, fill=stage6, below=12mm of s8] (s9) {\textbf{9.}~Frontend\\\footnotesize alertas + SHAP};

% Flechas track A
\draw[arrow] (s1) -- (s2);
\draw[arrow] (s2) -- (s2b);

% Flechas track B
\draw[arrow] (s3) -- (s4);
\draw[arrow] (s4) -- (s5);
\draw[arrow] (s5.south) |- (s6.east);
\draw[arrow] (s6) -- (s7);
\draw[arrow] (s7) -- (s8);
\draw[arrow] (s8) -- (s9);

% Cruce: featurize tambien sirve a inferencia
\draw[branch] (s4.south) to[bend right=20] (s8.north);

% Etiquetas tracks
\node[font=\small\bfseries, anchor=west, text=accent!70!black]
  at ([xshift=-3mm,yshift=2mm]s1.north west) {Track A --- Histórico / Seguimiento};
\node[font=\small\bfseries, anchor=west, text=accent!70!black]
  at ([xshift=-3mm,yshift=2mm]s3.north west) {Track B --- Predictivo};

\end{tikzpicture}
\end{center}

\noindent
Dos tracks paralelos comparten el mismo \texttt{featurize()}:
\begin{itemize}
  \item \textbf{Track A (Histórico)}: Excel real $\rightarrow$ SQLite $\rightarrow$ vista de Seguimiento (KPIs, SLA, motivos, anomalías). \emph{No usa el modelo.}
  \item \textbf{Track B (Predictivo)}: 60 días sintéticos $\rightarrow$ feature engineering $\rightarrow$ XGBoost calibrado $\rightarrow$ inferencia en vivo cada 3s sobre el plan del día.
\end{itemize}

% =====================================================================
\section{Stage 1 --- Datos crudos del Excel}
% =====================================================================

\begin{stagebox}{stage1}{Input: \texttt{datos\_eta\_2026\_04\_19.xlsx}\quad{}|\quad{}Output: 2 DataFrames pandas}
Hojas \texttt{Simpli} (2.120 filas) y \texttt{Geo} (2.221 filas).
Dump fiel del extracto operacional de SimpliRoute. \textbf{No se transforma nada en este stage}: solo se lee con \texttt{pd.read\_excel}.
\end{stagebox}

\subsection{Hoja \texttt{Simpli} --- columnas relevantes}

\begin{longtable}{p{4.4cm} p{2.3cm} p{8cm}}
\toprule
\textbf{Columna} & \textbf{Tipo} & \textbf{Rol en el pipeline} \\
\midrule
\endhead
\texttt{id} & BIGINT & PK natural; usada para dedupe en SQLite \\
\texttt{planned\_date} & DATE & Particionamiento + filtro temporal \\
\texttt{title} & TEXT & Razón social del cliente; matching de VIPs \\
\texttt{address} & TEXT & Dirección; mostrada en UI \\
\texttt{[order]} & INT & Posición en la ruta del vehículo \\
\texttt{checkout\_cl} & TIMESTAMP & Hora real de entrega \\
\texttt{current\_eta\_cl} & TIMESTAMP & ETA prometido \\
\texttt{status} & TEXT & \texttt{completed} / \texttt{failed} (label observado) \\
\texttt{sla\_hour\_checkout\_eta} & DECIMAL(10,4) & SLA en horas (signed); insumo de KPIs y bins \\
\texttt{bin\_label}, \texttt{bin\_index} & TEXT/INT & Bucket SLA pre-calculado para la distribución \\
\texttt{ct} & TEXT & Centro de transferencia (\texttt{CD OMNICANAL LOF2}, \dots) \\
\texttt{Patente\_falsa} & INT & Vehículo (anonimizado) \\
\texttt{Empresa\_falsa} & INT & \textbf{Tenant} --- multi-tenancy en POC y prod \\
\texttt{Drivername} & TEXT & Conductor \\
\texttt{Fechainicioruta} & TEXT & Inicio de ruta (string con timezone) \\
\texttt{am\_pm} & TEXT & \texttt{AM} / \texttt{PM} \\
\texttt{ruta\_eta\_futuro} & BIT & \multirow{4}{*}{\parbox{8cm}{4 flags + 1 agregada que el equipo de Falabella ya identificó como predictivas. Hoy se exponen en Seguimiento; \textbf{listas para ser features del modelo}.}} \\
\texttt{ruta\_fecha\_inicio\_mayor\_eta} & BIT & \\
\texttt{ruta\_primer\_punto\_lejano} & BIT & \\
\texttt{ruta\_fecha\_inicio\_distinta\_fecha\_eta} & BIT & \\
\texttt{ruta\_anomala} & BIT & \\
\bottomrule
\end{longtable}

\subsection{Hoja \texttt{Geo} --- columnas relevantes}

\begin{longtable}{p{3.6cm} p{2.3cm} p{8.6cm}}
\toprule
\textbf{Columna} & \textbf{Tipo} & \textbf{Rol en el pipeline} \\
\midrule
\endhead
\texttt{Suborden} & BIGINT & PK natural \\
\texttt{idruta} & BIGINT & Une con Simpli (vía \texttt{Fechainicioruta}) \\
\texttt{empresa\_falsa} & INT & Tenant (mismo eje que Simpli) \\
\texttt{localidad}, \texttt{region} & TEXT & Geografía $\rightarrow$ feature \texttt{comuna\_id} real (futuro) \\
\texttt{fechapactada} & DATE & Día comprometido al cliente \\
\texttt{tipodocumento} & TEXT & Boleta / factura / etc $\rightarrow$ feature posible \\
\texttt{estado} & TEXT & Estado terminal (\texttt{Entregado}, \texttt{Pendiente}, \dots) \\
\texttt{motivonoentrega} & TEXT & Razón cuando \texttt{estado} no es éxito $\rightarrow$ feature \\
\texttt{comentarionoentrega} & TEXT & Texto libre del driver \\
\bottomrule
\end{longtable}

% =====================================================================
\section{Stage 2 --- Carga a SQLite}
% =====================================================================

\begin{stagebox}{stage2}{Input: 2 DataFrames\quad|\quad{}Output: 6 tablas \texttt{fpoc\_*} en \texttt{valuedata.db}\quad|\quad{}Script: \texttt{fpoc\_loader/seed\_sqlite.py}}
Idempotente: borra y reinserta por \texttt{planned\_date} (Simpli) y por \texttt{idruta} (Geo). Ejecutar cuando llega un Excel nuevo.
\end{stagebox}

\subsection{Mapping de columnas Excel $\rightarrow$ tabla}

\begin{longtable}{p{4.5cm} p{4.5cm} p{6.2cm}}
\toprule
\textbf{Origen Excel} & \textbf{Tabla SQLite} & \textbf{Transformación} \\
\midrule
\endhead
Hoja Simpli (todas las cols) & \texttt{fpoc\_simpli\_visits} & Cast a tipos del schema; flags BIT $\rightarrow$ INTEGER 0/1; \texttt{Fechainicioruta\_hora\_cl} (TIME) $\rightarrow$ string \texttt{HH:MM:SS} \\
Hoja Geo (todas las cols) & \texttt{fpoc\_geo\_suborders} & \texttt{lpn}, \texttt{parentorder} a \texttt{Int64} (nullable); \texttt{NaN} $\rightarrow$ \texttt{NULL} \\
\texttt{DISTINCT Empresa\_falsa} & \texttt{fpoc\_empresas\_transporte} & 1 fila por empresa con nombre \texttt{Transporte NN} \\
(seed manual) & \texttt{fpoc\_users} & 1 admin + 1 ops + 1 transport\_manager por empresa, password bcrypt \\
\bottomrule
\end{longtable}

\subsection{Verificación post-carga}

\begin{lstlisting}[language=]
[verify] simpli=1866 geo=2119 anomalas=260 empresas=5 users={'falabella_admin': 1, 'falabella_ops': 1, 'transport_manager': 5}
\end{lstlisting}

% =====================================================================
\section{Stage 3 --- Generador sintético (insumo del modelo)}
% =====================================================================

\begin{stagebox}{stage3}{Input: SEED + parámetros\quad|\quad{}Output: 60 days $\times$ 120 visits $=$ $\sim$7.200 filas\quad|\quad{}Función: \texttt{pipeline.gen\_day\_visits()}}
Por qué sintético: el Excel actual es 1 solo día y no trae el evento histórico ``fallo anticipado vs ventana''. El generador reproduce los patrones operacionales (zonas problema, ventanas tardías, drivers conflictivos, recurrencia) para entrenar y validar el clasificador. \textbf{Reemplazable por un loader real de SimpliRoute en producción.}
\end{stagebox}

\subsection{Patrones ocultos (los que SHAP debe encontrar)}

\begin{itemize}
  \item \textbf{3 comunas problema}: celdas $0.05^{\circ} \times 0.05^{\circ}$ alrededor del depot $(-33.45, -70.66)$. Penalty $\times 3.0$.
  \item \textbf{2 drivers problema}: 2 vehículos de 12. Penalty $\times 2.5$.
\end{itemize}

\subsection{Receta del retraso}

\begin{lstlisting}[language=]
delay_i = local_noise * vehicle_factor * franja_factor(hora)
        + Exp(4) * (penalty_mult - 1)
        + 0.6 * delay_{i-1}
\end{lstlisting}

\textbf{Penalty multipliers acumulables}:

\begin{tabular}{lll}
\toprule
\textbf{Condición} & \textbf{Mult} & \textbf{Origen} \\
\midrule
Comuna problema & $\times 3.0$ & Hidden pattern \\
\texttt{window\_end} $\in \{17,18,19\}$ & $\times 2.0$ & Franja tarde \\
Driver problema & $\times 2.5$ & Hidden pattern \\
\texttt{load} $> 15$ & $\times 1.25$ & Carga \\
Cliente recurrente & $\times 4.0$ & Histórico cliente \\
\bottomrule
\end{tabular}

\subsection{Output del stage}

DataFrame con columnas: \texttt{id}, \texttt{tracking\_id}, \texttt{customer\_id}, \texttt{title}, \texttt{address}, \texttt{lat/lon}, \texttt{load}, \texttt{window\_start/end}, \texttt{planned\_arrival\_time}, \texttt{vehicle\_id}, \texttt{order}, \texttt{delay\_min}, \texttt{eta\_real}, \texttt{slack\_min}, \textbf{\texttt{failed}} (label).

% =====================================================================
\section{Stage 4 --- Feature engineering}
% =====================================================================

\begin{stagebox}{stage4}{Input: DataFrame con visitas + label\quad|\quad{}Output: matriz X ($\sim$70--90 cols) + vector y\quad|\quad{}Función: \texttt{pipeline.featurize()}}
Mismo \texttt{featurize()} se usa para entrenamiento e inferencia --- esto garantiza que las features que ve el modelo en producción son exactamente las que vio entrenando.
\end{stagebox}

\subsection{4.1 Features numéricos derivados}

\begin{longtable}{p{4.2cm} p{6cm} p{5cm}}
\toprule
\textbf{Feature} & \textbf{Cómputo} & \textbf{Hipótesis} \\
\midrule
\endhead
\texttt{hora\_window\_end} & \texttt{int(window\_end[:2])} & Tarde concentra fallos \\
\texttt{carga} & \texttt{load} (m$^3$/kg) & Carga grande $\rightarrow$ descarga lenta \\
\texttt{dist\_depot\_km} & Haversine cliente $\leftrightarrow$ depot & Más lejos $\rightarrow$ más viaje \\
\texttt{orden\_en\_ruta} & \texttt{order} (1\dots N) & Más tarde en ruta $\rightarrow$ más acum. \\
\texttt{retraso\_acumulado\_vehiculo} & \texttt{cumsum(delay\_min)} \emph{shifted} & Señal directa de ``ya viene tarde'' \\
\texttt{tasa\_fallo\_historica\_cliente} & Fail rate por \texttt{comuna\_id} en histórico & Codifica zona problema \\
\texttt{horas\_hasta\_window\_end} & \texttt{(window\_end $-$ ref\_clock)/3600}, clipped 0 & Anticipación: cuánto margen queda \\
\bottomrule
\end{longtable}

\subsection{4.2 Features categóricos (one-hot)}

\begin{tabular}{p{3.4cm} p{6cm} p{5cm}}
\toprule
\textbf{Feature} & \textbf{Cómputo} & \textbf{Cardinalidad típica} \\
\midrule
\texttt{comuna\_id} & Bucket $0.05^{\circ} \times 0.05^{\circ}$ de \texttt{(lat, lon)} & 50--80 zonas \\
\texttt{conductor\_id} & \texttt{v\{vehicle\_id\}} (ej.\ v1, v2, \dots) & 12 \\
\texttt{dia\_semana} & \texttt{planned\_date.dayofweek} & 7 \\
\bottomrule
\end{tabular}

\noindent
\texttt{pd.get\_dummies(prefix=["comuna","drv","dow"])} expande estas 3 a $\sim$70--100 columnas binarias. La lista exacta queda fija en \texttt{train\_cols} y se reaplica con \texttt{align\_columns} en inferencia.

\subsection{4.3 Anti-leak: \texttt{randomize\_observation}}

Sin este truco el modelo se haría trampa (vería el delay total cuando va a predecir antes del fin de ruta). Solución: para cada visita histórica se elige un ``reloj'' \texttt{\_ref} aleatorio entre 09:00 y \texttt{window\_end}, y todo lo que pasaría \emph{después} de ese reloj se setea a 0:

\begin{lstlisting}[language=pyplus]
span_sec   = (window_end - day_start).total_seconds()
offset_sec = span_sec * uniform(0, 1)
ref        = day_start + offset_sec

unobserved = eta_real > _ref
df.loc[unobserved, "delay_min"] = 0.0
retraso_acumulado = cumsum(delay_min) shifted   # solo lo ya completado
\end{lstlisting}

Esto reproduce el escenario real de inferencia (mirar a las 13:00 y predecir si la visita de las 18:00 va a fallar) \emph{durante} el entrenamiento.

\subsection{4.4 Resultado del stage}

\begin{itemize}
  \item Matriz \texttt{X\_full}: 7.200 filas $\times$ $\sim$80 columnas (mix numéricas + one-hot).
  \item Vector \texttt{y\_full}: 7.200 valores 0/1, base rate $\approx 20\%$.
  \item \texttt{train\_cols}: lista ordenada de columnas (se persiste para inferencia).
  \item \texttt{comuna\_rate}: dict \texttt{\{comuna\_id $\rightarrow$ fail\_rate\}} para enriquecer en inferencia.
\end{itemize}

% =====================================================================
\section{Stage 5 --- Split temporal}
% =====================================================================

\begin{stagebox}{stage5}{Criterio: por fecha (no random) para detectar drift entre días.}
\textbf{Train}: días 0--49 ($\approx$ 6.000 filas). \quad{}\textbf{Val}: días 50--59 ($\approx$ 1.200 filas).
\end{stagebox}

\begin{lstlisting}[language=pyplus]
dates_sorted = sorted(hist["planned_date"].unique())
train_dates  = set(dates_sorted[:50])
is_train     = hist_sorted["planned_date"].isin(train_dates).values

X_train = X_full.iloc[is_train]    # ~6000 filas
X_val   = X_full.iloc[~is_train]   # ~1200 filas
\end{lstlisting}

% =====================================================================
\section{Stage 6 --- Entrenamiento}
% =====================================================================

\begin{stagebox}{stage5}{Pipeline de entrenamiento: XGBoost $\rightarrow$ \texttt{CalibratedClassifierCV(method=isotonic)} $\rightarrow$ persistencia en \texttt{STATE.boot}.}
\end{stagebox}

\subsection{6.1 Modelo base (XGBoost)}

\begin{lstlisting}[language=pyplus]
spw = (y_train == 0).sum() / max(1, (y_train == 1).sum())   # ~4 (desbalance)

base_xgb = xgb.XGBClassifier(
    n_estimators=300,
    max_depth=5,
    learning_rate=0.05,
    scale_pos_weight=spw,
    eval_metric="logloss",
    tree_method="hist",
    n_jobs=-1, random_state=42,
)
\end{lstlisting}

\textbf{Por qué estos hiperparámetros}:

\begin{itemize}
  \item \texttt{n\_estimators=300, max\_depth=5, lr=0.05}: combinación clásica, robusta para datasets de $\sim$10k filas. Más profundidad sobreajustaría con tan poca data.
  \item \texttt{scale\_pos\_weight=spw}: pondera la clase minoritaria (fallos) para que el modelo no caiga en ``predecir siempre 0''.
  \item \texttt{tree\_method="hist"}: 3--5x más rápido que \texttt{exact}, sin pérdida significativa de calidad para este tamaño.
\end{itemize}

\subsection{6.2 Calibración isotonic}

\begin{lstlisting}[language=pyplus]
cal = CalibratedClassifierCV(base_xgb, method="isotonic", cv=3)
cal.fit(X_train.values, y_train)
\end{lstlisting}

XGBoost devuelve scores no calibrados (cercanos a $0$ o $1$ pero no probabilidades reales). \texttt{CalibratedClassifierCV} con \texttt{method="isotonic"} y 3-fold CV ajusta una función monótona que mapea \emph{score} $\rightarrow$ \emph{probabilidad real}, de modo que cuando el modelo dice ``70\%'' la frecuencia observada también es $\sim$70\%.

\subsection{6.3 SHAP (modelo paralelo, no calibrado)}

\texttt{TreeExplainer} de SHAP necesita acceso al booster crudo, que \texttt{CalibratedClassifierCV} oculta. Se entrena un segundo XGB \emph{idéntico} sobre el mismo set y se usa solo para explicar:

\begin{lstlisting}[language=pyplus]
shap_model = xgb.XGBClassifier(...).fit(X_train, y_train)
explainer  = shap.TreeExplainer(shap_model)
\end{lstlisting}

\textbf{Trade-off}: las probabilidades vienen del modelo calibrado; las atribuciones SHAP vienen del modelo no calibrado. Como ambos entrenan sobre la misma data y arquitectura, el ranking de features es estable; las contribuciones absolutas son aproximaciones.

% =====================================================================
\section{Stage 7 --- Pruebas y validación}
% =====================================================================

\begin{stagebox}{stage5}{Output: \texttt{STATE.boot["metrics"]} expuesto vía \texttt{/api/model/metrics}.}
\end{stagebox}

\subsection{7.1 Métricas}

\begin{tabular}{p{3.5cm} p{3cm} p{8.5cm}}
\toprule
\textbf{Métrica} & \textbf{Valor actual} & \textbf{Interpretación} \\
\midrule
ROC-AUC & $\approx 0.77$ & Capacidad discriminativa decente para 60d sintéticos. Esperable subir a 0.82--0.88 con 12 meses reales. \\
Brier score & $\approx 0.097$ & MSE de probabilidades: cercano al óptimo dado el base rate de $\approx 0.20$. \\
Confusion matrix (\texttt{thr}=0.50) & 2x2 & Permite calcular precision/recall/F1 con cualquier threshold. \\
Calibration curve & 10 bins & Confirma que ``predicho 0.7'' $\approx$ ``real 0.7'' (tras isotonic). \\
\bottomrule
\end{tabular}

\subsection{7.2 Cómo se valida en runtime}

\begin{itemize}
  \item Endpoint \texttt{GET /api/model/metrics} expone todo el dict.
  \item El frontend grafica la calibration curve (predicted vs actual) y el confusion matrix.
  \item Endpoint \texttt{GET /api/model/importance} entrega top-K features por \texttt{|SHAP|} promedio.
\end{itemize}

\subsection{7.3 Sanity checks que se ejecutan al boot}

\begin{enumerate}
  \item \textbf{Patrones ocultos visibles}: SHAP debe ranquear las features \texttt{drv\_v\{problema\}} y los \texttt{comuna\_*} de zonas problema entre el top-15.
  \item \textbf{Base rate train $\approx$ val}: $\approx 0.20$ en ambos $\rightarrow$ no hubo skew de muestreo.
  \item \textbf{Calibración monótona}: la curva no debe cruzarse a sí misma.
\end{enumerate}

% =====================================================================
\section{Stage 8 --- Inferencia en runtime}
% =====================================================================

\begin{stagebox}{stage6}{Función: \texttt{pipeline.apply\_status\_and\_predict()}\quad|\quad{}Disparado por: scheduler cada 3s + endpoints REST.}
\end{stagebox}

\subsection{8.1 Pasos por tick}

\begin{enumerate}
  \item Aplicar incidentes manuales/automáticos al \texttt{delay\_min} de visitas pendientes.
  \item Marcar \texttt{is\_completed = eta\_real $\leq$ sim\_clock}; calcular \texttt{\_obs\_delay} por vehículo (delay del último completed).
  \item Proyectar \texttt{current\_eta = planned\_arrival + \_obs\_delay} para pendientes.
  \item Computar \texttt{slack\_min} y \texttt{alert\_slack} (RED/YELLOW/GREEN).
  \item Llamar \texttt{featurize()} con \texttt{now\_clock=sim\_clock} (\textbf{misma función que en train}).
  \item \texttt{align\_columns(X, train\_cols)} $\rightarrow$ asegura mismo orden y rellena cols faltantes con 0.
  \item \texttt{cal\_model.predict\_proba(X)[:,1]} $\rightarrow$ \texttt{p\_fallo}.
  \item \texttt{explainer.shap\_values(X)} $\rightarrow$ matriz de atribuciones.
\end{enumerate}

\subsection{8.2 Reglas de alerta}

\begin{lstlisting}[language=]
alert_slack = "RED"    si slack_min <= 0
              "YELLOW" si  0 < slack_min <= 20
              "GREEN"  si slack_min > 20

alert_valuedata = (p_fallo >= 0.50)
                  AND (horas_hasta_window_end >= 2.0)
                  AND (status == "pending")
\end{lstlisting}

\texttt{alert\_valuedata} es la \emph{alerta anticipada} que diferencia la torre: dispara $\geq$ 2h antes del deadline cuando el modelo asigna $\geq$ 50\% de probabilidad de fallo. Las constantes son configurables y deberán tunearse contra ROI medido.

% =====================================================================
\section{Stage 9 --- Frontend (consumo)}
% =====================================================================

\begin{stagebox}{stage6}{Endpoints REST consumidos por React + Tailwind + Deck.gl.}
\end{stagebox}

\begin{tabular}{p{5.5cm} p{9.5cm}}
\toprule
\textbf{Endpoint} & \textbf{Qué entrega} \\
\midrule
\texttt{GET /api/state} & Reloj simulado, totales, vehículos, incidentes activos \\
\texttt{GET /api/kpis} & KPIs en vivo (completadas, alertas VD, rescate proyectado) \\
\texttt{GET /api/visits} & Listado completo con \texttt{p\_fallo}, \texttt{slack}, alertas \\
\texttt{GET /api/alerts/anticipated} & Top alertas VD ordenadas por \texttt{p\_fallo}, con top-3 SHAP \\
\texttt{GET /api/visits/\{tid\}/explanation} & Top-5 factores SHAP (signed) para una visita específica \\
\texttt{GET /api/model/metrics} & Calibration curve + confusion + AUC + Brier \\
\texttt{GET /api/model/importance} & Top-K features por \texttt{|SHAP|} promedio \\
\texttt{GET /api/watchlist} & Visitas pending scoreadas por urgencia \\
\texttt{GET /api/plan-diario} & Vista jerárquica empresa $\rightarrow$ vehículo $\rightarrow$ visitas \\
\bottomrule
\end{tabular}

% =====================================================================
\section{Camino a producción}
% =====================================================================

\begin{longtable}{p{3.7cm} p{3.6cm} p{4.6cm} p{2.6cm}}
\toprule
\textbf{Componente} & \textbf{POC} & \textbf{Producción} & \textbf{¿Reuso?} \\
\midrule
\endhead
Stage 1: Origen de datos & Excel del cliente & SimpliRoute REST API & Reemplazo \\
Stage 2: Persistencia & SQLite local & Azure SQL (schema \texttt{fpoc}) & Switch via \texttt{DB\_BACKEND} \\
Stage 3: Datos de train & Generador sintético 60d & Histórico real 12 meses + label observado & Reemplazo \\
Stage 4: Feature engineering & \texttt{featurize()} & \textbf{Idéntica} & \checkmark{} 100\% \\
Stage 5: Split & 50/10 días sintéticos & 11/1 meses reales & Adaptar \\
Stage 6: Modelo & XGBoost + isotonic + SHAP & \textbf{Idéntico} & \checkmark{} 100\% \\
Stage 7: Pruebas & AUC, Brier, calib & + monitoreo de drift mensual & Extender \\
Stage 8: Inferencia & cada 3s in-process & idem o batch programado & \checkmark{} 100\% \\
Stage 9: Frontend & React + Vite & Idéntico, build estático en SWA & \checkmark{} 100\% \\
\bottomrule
\end{longtable}

\textbf{Ganancias esperables con datos reales}: AUC 0.82--0.88 (+5--10 puntos), comuna real (no grid) +1--2 puntos, anomalías \texttt{ruta\_*} como features +1--3 puntos, contexto externo (clima, congestión, eventos) +1--3 puntos por feature.

% =====================================================================
\appendix
\section{Apéndice A --- Linaje completo columna por columna}
% =====================================================================

\begin{landscape}
\footnotesize
\begin{longtable}{p{4cm} p{3cm} p{3.5cm} p{4cm} p{4.5cm}}
\toprule
\textbf{Columna Excel / Sintética} & \textbf{Stage 2 (SQLite)} & \textbf{Stage 3 (Generador)} & \textbf{Stage 4 (Feature)} & \textbf{Uso final} \\
\midrule
\endhead
\texttt{Simpli.id} & \texttt{fpoc\_simpli\_visits.id} & --- & --- & PK, dedupe \\
\texttt{Simpli.planned\_date} & \texttt{planned\_date} (DATE) & --- & \texttt{dia\_semana} (one-hot) & Feature + filtro \\
\texttt{Simpli.title} & \texttt{title} & --- & --- & UI; matching VIP \\
\texttt{Simpli.address} & \texttt{address} & --- & --- & UI \\
\texttt{Simpli.[order]} & \texttt{"order"} & --- & \texttt{orden\_en\_ruta} & Feature \\
\texttt{Simpli.checkout\_cl} & \texttt{checkout\_cl} (TIMESTAMP) & análogo a \texttt{eta\_real} & --- (label \texttt{failed} se deriva) & KPIs Seguimiento \\
\texttt{Simpli.current\_eta\_cl} & \texttt{current\_eta\_cl} & análogo a \texttt{planned\_arrival\_time} & \texttt{horas\_hasta\_window\_end} & Feature \\
\texttt{Simpli.status} & \texttt{status} & \texttt{failed} (label) & Target \texttt{y} & Train target \\
\texttt{Simpli.sla\_hour\_checkout\_eta} & \texttt{sla\_hour\_checkout\_eta} & \texttt{slack\_min} (vía \texttt{eta\_real}) & --- & KPIs + alerta\_slack \\
\texttt{Simpli.bin\_label/index} & ídem & --- & --- & Distribución SLA \\
\texttt{Simpli.ct} & \texttt{ct} & --- & --- (futuro: feature) & Filtro UI \\
\texttt{Simpli.Patente\_falsa} & \texttt{Patente\_falsa} & \texttt{vehicle\_id} & \texttt{conductor\_id} (one-hot) & Feature \\
\texttt{Simpli.Empresa\_falsa} & \texttt{Empresa\_falsa} + \texttt{fpoc\_empresas\_transporte} & --- & --- & \textbf{Tenant} (multi-tenancy) \\
\texttt{Simpli.Drivername} & \texttt{Drivername} & --- (vehicle\_id sirve) & --- & UI \\
\texttt{Simpli.Fechainicioruta} & \texttt{Fechainicioruta} & --- & --- & Join con Geo \\
\texttt{Simpli.am\_pm} & \texttt{am\_pm} & --- & implícito en \texttt{hora\_window\_end} & Filtro UI \\
\texttt{Simpli.ruta\_*} (5 flags) & 5 columnas & --- & \textbf{listas para enchufar} & Hoy: Seguimiento. Futuro: features binarias \\
\midrule
\texttt{Geo.Suborden} & \texttt{Suborden} & --- & --- & PK, dedupe \\
\texttt{Geo.idruta} & \texttt{idruta} & --- & --- & Join con Simpli \\
\texttt{Geo.localidad/region} & ídem & --- & \textbf{candidato} a reemplazar grid de \texttt{comuna\_id} & Hoy: Seguimiento by-localidad \\
\texttt{Geo.tipodocumento} & ídem & --- & \textbf{candidato} feature one-hot & Hoy: nada \\
\texttt{Geo.estado} & ídem & --- & --- & Filtro/KPI \\
\texttt{Geo.motivonoentrega} & ídem & --- & \textbf{candidato} feature \texttt{motivo\_recurrente\_cliente} & Hoy: top motivos UI \\
\midrule
\multicolumn{5}{l}{\textit{Features generadas en Stage 4 que no tienen origen Excel directo (vienen del generador, derivadas en runtime):}} \\
\texttt{load} (sintético) & --- & sí, $U(0.5, 25)$ & \texttt{carga} & Feature; placeholder hasta tener carga real \\
\texttt{lat/lon} (sintético) & --- & sí, alrededor de DEPOT & \texttt{dist\_depot\_km}, \texttt{comuna\_id} & Feature \\
\texttt{delay\_min} (sintético) & --- & receta del Stage 3 & \texttt{retraso\_acumulado\_vehiculo} & Feature \\
\texttt{comuna\_id} (derivada) & --- & --- & sí (one-hot) & Feature \\
\texttt{tasa\_fallo\_historica\_cliente} (derivada) & --- & sí (\texttt{comuna\_rate}) & sí & Feature \\
\bottomrule
\end{longtable}
\end{landscape}

% =====================================================================
\section{Apéndice B --- Checklist de validación}
% =====================================================================

\textbf{Tras correr \texttt{seed\_sqlite.py}}:
\begin{itemize}[leftmargin=2em]
  \item[$\square$] \texttt{verify} reporta los counts esperados (\texttt{simpli}, \texttt{geo}, \texttt{empresas}, \texttt{users}).
  \item[$\square$] \texttt{SELECT COUNT(*) FROM fpoc\_simpli\_visits} $>$ 0.
  \item[$\square$] \texttt{SELECT COUNT(DISTINCT Empresa\_falsa)} = nro de filas en \texttt{fpoc\_empresas\_transporte}.
\end{itemize}

\textbf{Tras boot del backend} (\texttt{lifespan} en \texttt{main.py}):
\begin{itemize}[leftmargin=2em]
  \item[$\square$] Log: \texttt{Model ready. AUC=X.XXX, Brier=X.XXXX. Today plan: 120 visits.}
  \item[$\square$] \texttt{GET /api/model/metrics} devuelve dict completo.
  \item[$\square$] \texttt{GET /api/model/importance?top\_k=15} muestra al menos un \texttt{drv\_v*} y un \texttt{comuna\_*} en el top.
  \item[$\square$] AUC $\geq 0.70$ (umbral mínimo aceptable para POC).
\end{itemize}

\textbf{Tras login + plan en runtime}:
\begin{itemize}[leftmargin=2em]
  \item[$\square$] \texttt{GET /api/visits} devuelve 120 filas con \texttt{p\_fallo} entre 0 y 1.
  \item[$\square$] \texttt{GET /api/alerts/anticipated} devuelve $>$ 0 alertas con \texttt{horas\_hasta\_window\_end $\geq 2$}.
  \item[$\square$] \texttt{GET /api/visits/\{tid\}/explanation} devuelve 5 factores SHAP con nombres humanizados.
  \item[$\square$] Multi-tenant: login \texttt{transporte22@demo.cl} solo ve empresa 22 en \texttt{/api/empresas}.
\end{itemize}

\end{document}
